\documentclass[11pt]{article}
\usepackage[a4paper,margin=24mm]{geometry}
\usepackage{amsmath,amssymb}
\usepackage{xcolor}
\definecolor{ink}{HTML}{17344B}
\definecolor{accent}{HTML}{087F8C}
\setlength{\parindent}{0pt}
\setlength{\parskip}{10pt}
\begin{document}
\pagestyle{plain}
{\color{accent}\small FORMULA SEARCH / EXAMPLE DOCUMENT}

{\color{ink}\Huge Find the notation.\\[3pt] Keep the context.}

This short, original document is a playground for Formula Search. Open this PDF
in Chrome, open the extension, and try the queries printed below. These are
illustrative equations, not research findings.

\section*{1. A matrix score}
For compatible matrices, a scaled score can be written as
\begin{equation}
 S = \frac{QK^T}{\sqrt{d_k}}.
\end{equation}
Try the numerator with \emph{Contains this notation}. Superscripts, order,
and the denominator matter: the expression is not an unstructured bag of letters.

\section*{2. A simple objective}
A squared error with a penalty gives another expression to explore:
\begin{equation}
 L = (y-\hat{y})^2 + \lambda w^2.
\end{equation}
Try the squared weight. Copy the returned LaTeX, save a bookmark, or inspect
the original PDF excerpt from the result card.

\vfill
{\small\color{accent} Local search. Original PDF. Explicit matching modes.}
\newpage
{\color{accent}\small FORMULA SEARCH / EXAMPLE DOCUMENT}

{\color{ink}\Huge One symbol,\\[3pt] several contexts.}

\section*{3. Repeated notation}
The same squared weight can appear on a different page:
\begin{equation}
 R = \lambda w^2.
\end{equation}
Search the squared weight and use the page filter to keep only the occurrence
on this page. A bookmark identifies an occurrence, not only its formula text.

\section*{4. Scalar algebra is an explicit choice}
For commuting scalars, an expanded square is
\begin{equation}
 z = x^2 + 2xy + y^2.
\end{equation}
The algebraic matching mode has a different contract from notation search.
It is bounded, may report incomplete checks, and assumes commuting scalars.
Do not use that assumption for arbitrary matrix products.

\section*{What to expect}
This example exercises PDF text extraction. On supported arXiv documents,
Formula Search can also search TeX source and retain its file and line reference.
Some expressions cannot be mapped precisely to a PDF location; the interface
reports that uncertainty. Scanned equations require the optional OCR model
and a review before adding recognized notation to the session.

\vfill
{\small Created for the Formula Search project. See examples/README.md.}
\end{document}
